\documentclass[11pt]{article}

\usepackage[T1]{fontenc}
\usepackage[utf8]{inputenc}
\usepackage{lmodern}
\usepackage[margin=1in]{geometry}
\usepackage{amsmath,amssymb,amsthm,mathtools}
\usepackage{booktabs}
\usepackage{array}
\usepackage{enumitem}
\usepackage{microtype}
\usepackage{xcolor}
\usepackage{hyperref}

\hypersetup{
  colorlinks=true,
  linkcolor=blue!55!black,
  citecolor=blue!55!black,
  urlcolor=blue!55!black,
  pdftitle={Bivariate McKay Recurrences and Meromorphic Tensor-Curl Spectral Functions on Spherical Space Forms},
  pdfauthor={Luqman Mahmood}
}

\newtheorem{theorem}{Theorem}[section]
\newtheorem{proposition}[theorem]{Proposition}
\newtheorem{lemma}[theorem]{Lemma}
\newtheorem{corollary}[theorem]{Corollary}
\theoremstyle{definition}
\newtheorem{definition}[theorem]{Definition}
\newtheorem{convention}[theorem]{Convention}
\newtheorem{remark}[theorem]{Remark}

\DeclareMathOperator{\Sym}{Sym}
\DeclareMathOperator{\Hom}{Hom}
\DeclareMathOperator{\Irr}{Irr}
\DeclareMathOperator{\Res}{Res}
\newcommand{\TT}{\mathrm{TT}}
\newcommand{\Z}{\mathbb Z}
\newcommand{\C}{\mathbb C}
\newcommand{\e}{\mathrm e}
\newcommand{\tr}{\operatorname{tr}}
\newcommand{\id}{\mathrm{id}}
\newcommand{\one}{\mathbf 1}
\newcommand{\zH}{\zeta_{\mathrm H}}
\newcommand{\bplus}{b^{+}}
\newcommand{\bminus}{b^{-}}
\newcolumntype{L}[1]{>{\raggedright\arraybackslash}p{#1}}

\title{Bivariate McKay Recurrences and Meromorphic\\
Tensor-Curl Spectral Functions on Spherical Space Forms}
\author{Luqman Mahmood\thanks{Copyright 2026 Luqman Mahmood. This paper is licensed under Creative Commons Attribution 4.0 International (CC BY 4.0); see \url{https://creativecommons.org/licenses/by/4.0/}. Software is excluded from this grant.}}
\date{1 October 2026}

\begin{document}

\maketitle

\begin{abstract}
Let $\Gamma$ be a finite subgroup of
$\operatorname{Spin}(4)=SU(2)_L\times SU(2)_R$, let
$V_n=\Sym^n\C^2$, and let $\boldsymbol m_{p,q}$ record the irreducible
$\Gamma$-multiplicities in $V_p\boxtimes V_q$.  If $A_L$ and $A_R$
denote tensor multiplication by the restricted left and right defining
representations, we prove
\[
 \boldsymbol m_{p,q}
 =U_p(A_L/2)U_q(A_R/2)\boldsymbol e_0,
 \qquad [A_L,A_R]=0.
\]
Thus every finite $\Gamma$ carries a bivariate representation-ring
recurrence, without a product assumption.  Restriction to the two
transverse-traceless helicity diagonals gives distinct rank operators
$A_L^2-2I$ and $A_R^2-2I$ and a diagonal rank--level identity.  For a
free action on $S^3$, these vectors are the twisted multiplicities of the
normalized rank-$J$ tensor-curl operator.

Classwise multiplicities are quasipolynomial in the level, so each branch
spectral function is a finite Hurwitz-zeta combination, meromorphic on
$\C$ with possible poles only at $1,2,3$.  The rank recurrence propagates
for all complex $s$ at fixed spectral offset; at the geometric offset
$c_J=J+1$, unilateral reindexing contributes an explicit lowest-mode
boundary term.  The signed branch difference is entire for every finite
subgroup.  For a free effective action, the remaining poles are central
and satisfy $\Res_2=0$ and $\Res_1=-J^2\Res_3$, yielding a convergent
small-time heat expansion and exact relative log-determinant formulas.
The cumulative identity
$\sum_{r=1}^{J}U_{2r-1}=U_{J-1}U_J$ supplies the corresponding finite
polynomial compression used in companion eta-transfer and boundary
constructions.  The results are confined to round spherical quotients;
the nonround problem remains open.
\end{abstract}

\section{Introduction}

The spectrum of a round spherical quotient is controlled by two finite
pieces of information: the $SU(2)_L\times SU(2)_R$ harmonic content of
$S^3$ and the restriction of those representations to the deck group.
The classical spectral literature turns this observation into character
formulas and one-variable generating functions
\cite{Ikeda1980I,Ikeda1980II,Ikeda1980Lens,Bar1996,
LauretMiatelloRossetti2021}.  The classical McKay correspondence organizes
restriction along one $SU(2)$ factor by an affine-ADE adjacency matrix
\cite{McKay1980,Kostant1985,Rossmann2003,Suter2007,
Benkart2016,JingWangZhang2020,JingWangZhang2021}.
For finite subgroups of $SO(4)$, Xiong introduced the corresponding
left/right pair as an edge-coloured McKay graph \cite{Xiong2020}.

This paper identifies the spectral recurrence carried by that pair.  The
basic object is not initially an eta value but the full bivariate
multiplicity grid
\[
 (p,q)\longmapsto
 \boldsymbol m_{p,q}
 =
 \left(
 \dim\Hom_\Gamma(\rho_i,V_p\boxtimes V_q)
 \right)_{\rho_i\in\Irr(\Gamma)} .
\]
The two Clebsch--Gordan rules act independently on this grid.  Their
commutativity produces a two-variable Chebyshev reconstruction from a
single source at $(0,0)$.  The tensor-curl spectrum is then obtained by
restricting to the two diagonals $q-p=\pm2J$.
Explicitly,
\[
 \boldsymbol m_{p,q}
 =U_p(A_L/2)U_q(A_R/2)\boldsymbol e_0 .
\]
This formula places the one-sided affine-ADE recurrences and genuinely
two-sided spherical space forms in the same finite algebraic system.

Three distinctions are essential.

\begin{enumerate}[label=(\roman*)]
\item The universal theorem is a representation-ring theorem for a finite
subgroup of $\operatorname{Spin}(4)$.  Literal affine-ADE graph language
applies to one-sided factors and product reductions, not to an arbitrary
two-sided subgroup.
\item The positive and negative helicity lines generally propagate with
different matrices.  Their signed and unsigned assemblies therefore obey
a coupled two-matrix identity.  The branch laws do not furnish a common
single-matrix recurrence, although special invariant-subspace
coincidences are not excluded.
\item A recurrence of multiplicities does not automatically become the
same recurrence for the geometric spectral functions at successive ranks.
The geometric eigenvalue is $n+J+1$, so its origin moves with $J$.  The
resulting correction is a unilateral boundary term, not a
group-cohomological construction.
\end{enumerate}

The main statements are as follows.  Theorem~\ref{thm:bivariate} gives the
two-variable representation-ring recurrence, commutation, Chebyshev
solution, and rational generating function.  Corollary~\ref{cor:cumulative-odd-chebyshev}
records the generic finite cumulative odd-Chebyshev identity, while
Remark~\ref{rem:p3-cumulative-interface} explains its use in the companion
eta-transfer and boundary constructions
\cite{Mahmood2026DiracToTensorCurlEtaTransfer,
Mahmood2026SteinWeissAPSBoundary}.
Theorem~\ref{thm:tt-lines}
restricts it to the tensor-curl helicity lines and proves the diagonal
rank--level identity.  Proposition~\ref{prop:finite-hurwitz} gives the
finite-Hurwitz continuation of every branch.  Theorem~\ref{thm:all-s}
proves all-$s$ propagation at fixed offset, while
Theorem~\ref{thm:physical-offset} gives the exact geometric-offset boundary
correction.  Corollary~\ref{cor:relative-det} records the relative zeta and
log-determinant consequences.  Theorem~\ref{thm:eta-entire} proves
eta-entirety and identifies all possible pole-producing conjugacy
classes, while Proposition~\ref{prop:heat-expansion} converts those
residues into a convergent heat expansion.  Finally,
Propositions~\ref{prop:relative-pole-free} and
\ref{prop:closed-logdet} give pole-free relative rows and their exact
Hurwitz-derivative form.

\paragraph{Scope.}
The finite-group representation theory, Clebsch--Gordan identity,
one-sided McKay correspondence, character projection, Hurwitz
continuation, and round tensor spectra are antecedent machinery.  The
contribution here is their combination into a bivariate recurrence and its
round spectral consequences, rather than a new regularization method.  The
cumulative Chebyshev identity is purely algebraic.  The Dirac-to-tensor-curl
eta transfer and its combination with the auxiliary eta, kernel, and
Atiyah--Patodi--Singer boundary formulas are established in the companion
papers cited above.  Extending the present spectral propagation to nonround
metrics remains open.

\section{Representation and spectral conventions}

\subsection{The bivariate restriction grid}

Let
\[
 \Gamma< SU(2)_L\times SU(2)_R
\]
be finite.  Write
\[
 \Irr(\Gamma)=\{\rho_0,\rho_1,\ldots,\rho_r\},
 \qquad \rho_0=\one,
 \qquad d_i=\dim\rho_i,
\]
and let $\boldsymbol e_0$ be the coordinate vector of the trivial
representation.  For $n\geq0$, put $V_n=\Sym^n\C^2$; thus
$\dim V_n=n+1$ and $V_n$ is self-dual.  Set
\[
 W_L=(V_1\boxtimes V_0)|_\Gamma,
 \qquad
 W_R=(V_0\boxtimes V_1)|_\Gamma .
\]
Our McKay-matrix convention is
\begin{equation}
 (A_L)_{ij}
 =\dim\Hom_\Gamma(\rho_i,W_L\otimes\rho_j),
 \qquad
 (A_R)_{ij}
 =\dim\Hom_\Gamma(\rho_i,W_R\otimes\rho_j).
 \label{eq:McKay-matrices}
\end{equation}
Because $W_L$ and $W_R$ are self-dual, both matrices are symmetric
nonnegative integer matrices.  Define
\begin{equation}
 m_{p,q}^i
 =\dim\Hom_\Gamma(\rho_i,(V_p\boxtimes V_q)|_\Gamma),
 \qquad
 \boldsymbol m_{p,q}=(m_{p,q}^i)_{i=0}^r ,
 \label{eq:multiplicity-grid}
\end{equation}
and use the zero-extension convention
$\boldsymbol m_{p,q}=0$ if $p<0$ or $q<0$.

\subsection{The round tensor-curl branches}

For the unit round $S^3$, the normalized projected
transverse-traceless tensor-curl operator of integer rank $J\geq1$ has
the two $SU(2)_L\times SU(2)_R$ branches
\begin{equation}
 \mathcal H_{J,n}^{+}=V_{n+2J}\boxtimes V_n,
 \qquad
 \mathcal H_{J,n}^{-}=V_n\boxtimes V_{n+2J},
 \qquad n\geq0,
 \label{eq:TT-branches}
\end{equation}
with eigenvalues $+(n+J+1)$ and $-(n+J+1)$ in the orientation used
here.  Reversing the orientation exchanges the branches and reverses
the signed eta function.  The local self-gradient construction and its
TT compression are standard higher-spin data; we use only the
representation-theoretic spectral realization
\cite{Homma2001,HommaTomihisa2021}.

We write $A_J$ for the unit-radius self-adjoint realization determined
by this orthogonal spectral decomposition on the $L^2$ space of TT
rank-$J$ tensors.  Its domain consists of the expansions for which the
sum of eigenvalue squared times coefficient norm squared is finite.
The displayed spectrum makes its resolvent compact.  On a free quotient
and in a twist sector, $A_J$ means the corresponding reducing
isotypic restriction.  In the absolute zeta function below, the
negative branch contributes the magnitudes $n+J+1$; no spectral-cut
phase is included.

If the induced action
\[
 (g_L,g_R)\cdot x=g_Lxg_R^{-1},\qquad x\in SU(2)\cong S^3,
\]
is effective and free, the quotient is a round spherical space form.
For each $\rho_i$, the copy-multiplicity convention for the associated
flat bundle gives
\begin{equation}
 \boldsymbol b_J^+(n)=\boldsymbol m_{n+2J,n},
 \qquad
 \boldsymbol b_J^-(n)=\boldsymbol m_{n,n+2J}.
 \label{eq:branch-multiplicities}
\end{equation}
The algebraic results below hold for every finite $\Gamma$.  Freeness is
needed only for this manifold interpretation.  An element
$(g_L,g_R)\notin\{(1,1),(-1,-1)\}$ has a fixed point precisely when
$g_L$ and $g_R$ are conjugate in $SU(2)$, equivalently when their
defining traces agree.

\begin{remark}[Central kernel and descent]
\label{rem:central-descent}
If $(-1,-1)\in\Gamma$, the action on $S^3$ is ineffective.  A central
element $z=(\varepsilon_L,\varepsilon_R)$ acts on
$V_p\boxtimes V_q$ by
$\varepsilon_L^p\varepsilon_R^q$.  Schur's lemma therefore forces
$m_{p,q}^i=0$ unless
\[
 \rho_i(z)=\varepsilon_L^p\varepsilon_R^q I_{\rho_i}.
\]
Every TT line has even total parity $p+q=2n+2J$ and hence descends as
an ordinary representation of the effective quotient by
$\langle(-1,-1)\rangle$.  Odd-parity spinorial lines instead retain
the chosen spin-lift data.  No spinorial transfer is used below.
\end{remark}

\section{The bivariate representation-ring theorem}

\begin{theorem}[Bivariate McKay recurrence]
\label{thm:bivariate}
For all $p,q\geq0$,
\begin{align}
 \boldsymbol m_{p+1,q}+\boldsymbol m_{p-1,q}
 &=A_L\boldsymbol m_{p,q},
 \label{eq:left-rec}\\
 \boldsymbol m_{p,q+1}+\boldsymbol m_{p,q-1}
 &=A_R\boldsymbol m_{p,q}.
 \label{eq:right-rec}
\end{align}
Moreover,
\begin{equation}
 [A_L,A_R]=0
 \label{eq:commuting}
\end{equation}
and the unique quadrant solution with source
$\boldsymbol m_{0,0}=\boldsymbol e_0$ is
\begin{equation}
 \boxed{\;
 \boldsymbol m_{p,q}
 =U_p(A_L/2)U_q(A_R/2)\boldsymbol e_0 ,
 \;}
 \label{eq:Chebyshev-solution}
\end{equation}
where $U_n$ is the Chebyshev polynomial of the second kind.  Equivalently,
\begin{equation}
 \sum_{p,q\geq0}\boldsymbol m_{p,q}t^pu^q
 =
 (I-tA_L+t^2I)^{-1}
 (I-uA_R+u^2I)^{-1}\boldsymbol e_0 .
 \label{eq:bivariate-GF}
\end{equation}
\end{theorem}

\begin{proof}
The Clebsch--Gordan identity
\[
 V_1\otimes V_p\cong V_{p+1}\oplus V_{p-1},
 \qquad V_{-1}=0,
\]
holds also at $p=0$.  Restricting to $\Gamma$, tensoring
$(V_p\boxtimes V_q)|_\Gamma$ by $W_L$, and projecting onto $\rho_i$
gives \eqref{eq:left-rec} by the definition
\eqref{eq:McKay-matrices}.  The right identity is identical with the
two factors exchanged.

The matrices are multiplication by $[W_L]$ and $[W_R]$ in the
commutative complex representation ring $R(\Gamma)$, proving
\eqref{eq:commuting}.  Since
\[
 U_{n+1}(x)=2xU_n(x)-U_{n-1}(x),
 \qquad U_{-1}=0,\quad U_0=1,
\]
the right-hand side of \eqref{eq:Chebyshev-solution} satisfies both
recurrences and the quadrant boundary condition.  The two three-term
recurrences determine the grid uniquely from $\boldsymbol e_0$, which
proves \eqref{eq:Chebyshev-solution}.  Summing each recurrence, or using
the standard generating series for $U_n$, gives
\eqref{eq:bivariate-GF}.
\end{proof}

\begin{corollary}[Cumulative odd-Chebyshev identity]
\label{cor:cumulative-odd-chebyshev}
For every integer $J\geq1$,
\begin{equation}
 \sum_{r=1}^{J}U_{2r-1}(x)=U_{J-1}(x)U_J(x)
 \label{eq:cumulative-odd-chebyshev}
\end{equation}
as an identity of polynomials in $x$.  Consequently, the left and right
representation-ring generators satisfy, separately,
\begin{align}
 \sum_{r=1}^{J}U_{2r-1}(A_L/2)
 &=U_{J-1}(A_L/2)U_J(A_L/2),
 \label{eq:cumulative-odd-chebyshev-left}\\
 \sum_{r=1}^{J}U_{2r-1}(A_R/2)
 &=U_{J-1}(A_R/2)U_J(A_R/2).
 \label{eq:cumulative-odd-chebyshev-right}
\end{align}
\end{corollary}

\begin{proof}
For $x=\cos\theta$ with $\sin\theta\ne0$, the formula
$U_n(\cos\theta)=\sin((n+1)\theta)/\sin\theta$ gives
\[
 U_JU_{J+1}-U_{J-1}U_J
 =\frac{\sin((J+1)\theta)
 [\sin((J+2)\theta)-\sin(J\theta)]}{\sin^2\theta}
 =U_{2J+1}.
\]
Both sides are polynomials, so this identity holds for every $x$.  The
case $J=1$ of \eqref{eq:cumulative-odd-chebyshev} is $U_1=U_0U_1$.
If it holds at $J$, the displayed product identity gives
\[
 \sum_{r=1}^{J+1}U_{2r-1}
 =U_{J-1}U_J+U_{2J+1}=U_JU_{J+1},
\]
which proves \eqref{eq:cumulative-odd-chebyshev} by induction.  Evaluating
this polynomial identity at $A_L/2$ and at $A_R/2$ proves the two matrix
specializations.
\end{proof}

\begin{remark}[Application to cumulative boundary eta]
\label{rem:p3-cumulative-interface}
The transfer theorem of
\cite{Mahmood2026DiracToTensorCurlEtaTransfer} denotes the symmetric-power
multiplication matrices by $A_m^X$ for $X=L,R$.  Under the compatible
representation-ring conventions, its Clebsch--Gordan recurrence gives
$A_m^X=U_m(A_X/2)$, so Corollary~\ref{cor:cumulative-odd-chebyshev}
simplifies the finite odd-rank sums in the cumulative assembly of
\cite{Mahmood2026SteinWeissAPSBoundary}.  Combining the rankwise
transfer with that paper's auxiliary eta, kernel, and
Atiyah--Patodi--Singer assembly yields the cumulative boundary
eta and reduced correction in the latter paper.  The present corollary
supplies the polynomial identity; the transfer and boundary theorems retain
their respective geometric hypotheses.
\end{remark}

\begin{remark}[Relation to coloured McKay graphs]
The pair $(A_L,A_R)$ itself is not new: in coloured-graph form it is
the finite-$SO(4)$ construction studied by Xiong \cite{Xiong2020}.
Theorem~\ref{thm:bivariate} adds the copy-multiplicity grid, the two
commuting recurrences, their Chebyshev reconstruction, and the spectral
restriction below.  For arbitrary $\Gamma$, the safest name is a
two-sided representation-ring McKay recurrence.  An arbitrary
two-coloured graph here need not itself be a single affine-ADE diagram.
\end{remark}

\section{Tensor-curl helicity lines}

Put
\begin{equation}
 M_+=A_L^2-2I,
 \qquad
 M_-=A_R^2-2I.
 \label{eq:rank-matrices}
\end{equation}
For algebraic boundary formulas, we also allow the formal seed $J=0$ in
\eqref{eq:branch-multiplicities}; the geometric tensor-curl family
begins at $J=1$.

\begin{theorem}[Rank and diagonal rank--level laws]
\label{thm:tt-lines}
For every $J\geq1$ and $n\geq0$,
\begin{align}
 \boldsymbol b_{J+1}^{+}(n)+\boldsymbol b_{J-1}^{+}(n)
 &=M_+\boldsymbol b_J^{+}(n),
 \label{eq:plus-rank}\\
 \boldsymbol b_{J+1}^{-}(n)+\boldsymbol b_{J-1}^{-}(n)
 &=M_-\boldsymbol b_J^{-}(n).
 \label{eq:minus-rank}
\end{align}
With $\boldsymbol b_J^\pm(-1)=0$, the positive branch also satisfies
\begin{align}
 A_LA_R\boldsymbol b_J^+(n)
 ={}&\boldsymbol b_J^+(n+1)+\boldsymbol b_J^+(n-1)
 \nonumber\\
 &+\boldsymbol b_{J-1}^+(n+1)
 +\boldsymbol b_{J+1}^+(n-1),
 \label{eq:diagonal-plus}
\end{align}
and the negative branch satisfies the mirror identity obtained by
interchanging $L$ and $R$.
\end{theorem}

\begin{proof}
Apply \eqref{eq:left-rec} twice at
$(p,q)=(n+2J,n)$:
\[
 A_L^2\boldsymbol m_{n+2J,n}
 =
 \boldsymbol m_{n+2J+2,n}
 +2\boldsymbol m_{n+2J,n}
 +\boldsymbol m_{n+2J-2,n}.
\]
This is \eqref{eq:plus-rank}; the right-factor calculation gives
\eqref{eq:minus-rank}.

Applying $A_LA_R$ to $\boldsymbol m_{n+2J,n}$ produces the four
corners
\[
 \boldsymbol m_{n+2J\pm1,n\pm1}.
\]
Reading those corners on the diagonal lines yields
\eqref{eq:diagonal-plus}.  Commutation makes the order of the two
steps immaterial, and the mirror calculation proves the negative
identity.
\end{proof}

\begin{remark}[Coupled total identities]
\label{rem:no-common}
Put $\boldsymbol u_J=\boldsymbol b_J^++\boldsymbol b_J^-$ and
$\boldsymbol v_J=\boldsymbol b_J^+-\boldsymbol b_J^-$.  The branch
laws give the exact coupled system
\begin{align*}
 \boldsymbol u_{J+1}+\boldsymbol u_{J-1}
 &=
 \tfrac12(M_++M_-)\boldsymbol u_J
 +\tfrac12(M_+-M_-)\boldsymbol v_J,\\
 \boldsymbol v_{J+1}+\boldsymbol v_{J-1}
 &=
 \tfrac12(M_+-M_-)\boldsymbol u_J
 +\tfrac12(M_++M_-)\boldsymbol v_J.
\end{align*}
Thus the theorem supplies two branch matrices rather than one matrix
closing on either total alone.  When $M_+\ne M_-$ and both
$\boldsymbol u_J$ and $\boldsymbol v_J$ have nonzero components, the
coupling terms are visible in the displayed system, so the
representation-ring identities alone do not yield a universal
common-matrix recurrence.  A special solution may nevertheless lie in a
common invariant subspace.  Spectral propagation below is therefore stated
branchwise.
\end{remark}

\section{Quasipolynomials and finite-Hurwitz continuation}

\begin{lemma}[Classwise quasipolynomiality]
\label{lem:quasipolynomial}
For fixed $J$, sign $\varepsilon\in\{+,-\}$, and irrep $\rho_i$, the
sequence $n\mapsto b_{J,i}^{\varepsilon}(n)$ is a quasipolynomial of
degree at most two.  A valid period is the least common multiple of
$2$ and the orders of the noncentral eigenphases occurring in the two
defining slots of the conjugacy classes of $\Gamma$.
\end{lemma}

\begin{proof}
Character projection gives
\begin{equation}
 m_{p,q}^i
 =\frac1{|\Gamma|}
 \sum_{g=(g_L,g_R)\in\Gamma}
 \overline{\chi_i(g)}\,
 \chi_p(g_L)\chi_q(g_R).
 \label{eq:character-projection}
\end{equation}
If $g_L$ has eigenvalues $\e^{\pm i\theta}$ with
$\theta\notin\{0,\pi\}$, then
\[
 \chi_p(g_L)=\frac{\sin((p+1)\theta)}{\sin\theta}
\]
is periodic in $p$ because $g_L$ has finite order.  At $g_L=\pm I$,
\[
 \chi_p(\pm I)=(\pm1)^p(p+1),
\]
a period-two sign times a linear polynomial.  The same alternatives
hold in the right slot.  On $p=n$, $q=n+2J$, or with the slots
exchanged, the product is polynomial of degree at most two on every
common residue class.  The finite sum
\eqref{eq:character-projection} preserves the claim.
\end{proof}

\begin{definition}[Branch spectral function]
\label{def:branch-F}
For $c>0$ and $\Re s>3$, define the vector-valued Dirichlet series
\begin{equation}
 \boldsymbol F_J^\pm(s;c)
 =\sum_{n=0}^{\infty}
 \boldsymbol b_J^\pm(n)(n+c)^{-s}.
 \label{eq:branch-F}
\end{equation}
The geometric offset is $c_J=J+1$.  Componentwise, the signed eta and
unsigned absolute spectral zeta functions are
\begin{align}
 \eta_{J,i}(s)
 &=F_{J,i}^+(s;c_J)-F_{J,i}^-(s;c_J),
 \label{eq:eta-from-branches}\\
 \zeta_{|A_J|,i}(s)
 &=F_{J,i}^+(s;c_J)+F_{J,i}^-(s;c_J).
 \label{eq:zeta-from-branches}
\end{align}
\end{definition}

\begin{proposition}[Finite-Hurwitz formula]
\label{prop:finite-hurwitz}
Every component of \eqref{eq:branch-F} extends meromorphically to
$\C$ as a finite linear combination of Hurwitz zeta functions.  Its
only possible poles are simple poles at $s=1,2,3$.
\end{proposition}

\begin{proof}
Let $a(n)$ be one component and choose a quasipolynomial period $P$.
For $0\leq r<P$, write on that residue class
\[
 a(P\ell+r)
 =\sum_{h=0}^{2}q_{r,h}(P\ell+r+c)^h .
\]
This is merely a change of polynomial basis and may have vanishing
top coefficients.  Then
\begin{equation}
 \sum_{n\geq0}a(n)(n+c)^{-s}
 =
 \sum_{r=0}^{P-1}\sum_{h=0}^{2}
 q_{r,h}P^{h-s}
 \zH\!\left(s-h,\frac{r+c}{P}\right).
 \label{eq:finite-Hurwitz}
\end{equation}
The Hurwitz zeta function is meromorphic with its only pole at first
argument $1$, proving the assertion.
\end{proof}

\begin{remark}[Terminology]
Some lens-space literature uses ``bivariate representations'' for
one-group highest weights of the form
$k\varepsilon_1+\ell\varepsilon_2$
\cite{LauretRossiBertone2018,LauretMiatelloRossetti2021}.
Here \emph{bivariate} instead
refers to the two independent symmetric-power indices $(p,q)$ of
$SU(2)_L\times SU(2)_R$.
\end{remark}

\section{Pole localization, eta entirety, and heat coefficients}

For a conjugacy-class representative $g=(g_L,g_R)$, choose folded
eigenangles $\theta_L,\theta_R\in[0,\pi]$ so that the defining
eigenvalues are $\e^{\pm i\theta_L}$ and
$\e^{\pm i\theta_R}$.  A class is \emph{central--central} when both
angles lie in $\{0,\pi\}$.  Write the corresponding signs as
$\varepsilon_L,\varepsilon_R\in\{\pm1\}$ and put
$\sigma=\varepsilon_L\varepsilon_R$.

\begin{theorem}[Eta entirety and central pole localization]
\label{thm:eta-entire}
Let $\Gamma< SU(2)_L\times SU(2)_R$ be finite, let $J\geq1$, and let
$c>0$.  For every irreducible coordinate $i$, the signed function
\begin{equation}
 E_{J,i}(s;c)=F^+_{J,i}(s;c)-F^-_{J,i}(s;c)
 \label{eq:eta-entire-function}
\end{equation}
is entire.

If the induced action on $S^3$ is free, all poles of either branch
come from central--central classes with $\sigma=+1$; for a free effective
action the identity is the only such class.  At the geometric
offset $c=J+1$, each branch has
\begin{align}
 \Res_{s=3}F^{\pm}_{J,i}(s;J+1)
 &=\frac1{|\Gamma|}
   \sum_{\substack{C\ \mathrm{central--central}\\\sigma(C)=+1}}
   |C|\,\overline{\chi_i(C)},
 \label{eq:residue-three}\\
 \Res_{s=2}F^{\pm}_{J,i}(s;J+1)&=0,
 \label{eq:residue-two}\\
 \Res_{s=1}F^{\pm}_{J,i}(s;J+1)
 &=-J^2\Res_{s=3}F^{\pm}_{J,i}(s;J+1).
 \label{eq:residue-one}
\end{align}
If the action is not free, a noncentral resonant class with
$\theta_L=\theta_R=\theta$ contributes additionally
\begin{equation}
 \frac{|C|\,\overline{\chi_i(C)}}{|\Gamma|}
 \frac{\cos(2J\theta)}{2\sin^2\theta}
 \label{eq:resonant-residue}
\end{equation}
to the residue at $s=1$ of \emph{both} branches.  Thus such resonances
do not obstruct the entirety of \eqref{eq:eta-entire-function}.
\end{theorem}

\begin{proof}
Character projection reduces the statement to one conjugacy class at
a time.  On a central--central class, the two branch factors are
identical:
\[
 \varepsilon_L^{n+2J}\varepsilon_R^n(n+2J+1)(n+1)
 =
 \varepsilon_L^n\varepsilon_R^{n+2J}(n+1)(n+2J+1),
\]
because the exponent $2J$ is even.  If $\sigma=-1$, this is an
alternating polynomial and every residue-class coefficient mean
vanishes.  If $\sigma=+1$, then at $c=J+1$ its polynomial part is
\[
 (n+1)(n+2J+1)=(n+c)^2-J^2,
\]
which gives \eqref{eq:residue-three}--\eqref{eq:residue-one}.

If exactly one slot is central, the remaining sine quotient is a
mean-zero periodic sequence; multiplication by the at-most-linear
central factor therefore produces no residue.  If neither slot is
central, the class sequence is bounded, so only a pole at $s=1$ is
possible.  Product-to-sum gives, for the plus branch, a difference of
cosines with frequencies $\theta_L-\theta_R$ and
$\theta_L+\theta_R$.  Folded noncentral angles make the second
frequency nonresonant.  The first has nonzero mean precisely when
$\theta_L=\theta_R$, and its mean is the quantity in
\eqref{eq:resonant-residue}.  Swapping the slots gives the same mean
on that locus.  A noncentral equality of folded angles is exactly the
fixed-vector obstruction to freeness.  Hence the branch residue maps
are equal for every finite group, proving that their difference is
entire, while freeness leaves only the stated central terms.
\end{proof}

The theorem is stronger than regularity of eta merely at the origin:
the entire meromorphic pole part cancels between the two helicities.
The unsigned function retains local-density information in its
central residues.

\begin{proposition}[Convergent heat expansion]
\label{prop:heat-expansion}
Let $F(s)=\sum_{n\geq0}a(n)(n+c)^{-s}$ be any component of a branch
function, let $P$ be a quasipolynomial period, and write
$r_j=\Res_{s=j}F(s)$ for $j=1,2,3$.  Then the first-order heat trace
\[
 \Theta_F(\tau)=\sum_{n\geq0}a(n)\e^{-(n+c)\tau}
\]
has the convergent Laurent expansion
\begin{equation}
 \Theta_F(\tau)
 =\sum_{j=1}^{3}\Gamma(j)r_j\tau^{-j}
  +\sum_{k=0}^{\infty}\frac{(-1)^k}{k!}F(-k)\tau^k,
 \qquad |\tau|<\frac{2\pi}{P}.
 \label{eq:heat-expansion}
\end{equation}
All coefficients are finite exact combinations of rational numbers
and Bernoulli-polynomial values.
\end{proposition}

\begin{proof}
Split the quasipolynomial into its $P$ residue classes.  Each class is
a degree-at-most-two polynomial times a geometric series in
$\e^{-P\tau}$, hence a rational function whose nearest nonzero
singularities occur at $P\tau=\pm2\pi i$.  Its Laurent coefficients
are obtained from
$\zH(-m,x)=-B_{m+1}(x)/(m+1)$.  Summing the finitely many classes
gives \eqref{eq:heat-expansion}; equivalently, the negative powers are
the residues of $\Gamma(s)F(s)\tau^{-s}$ and the nonnegative powers
are its residues at the nonpositive integers.
\end{proof}

\section{Meromorphic rank propagation}

\begin{theorem}[Fixed-offset all-$s$ identity]
\label{thm:all-s}
Fix $c>0$.  For every $J\geq1$, the following identities hold as
identities of vector-valued meromorphic functions on $\C$:
\begin{align}
 \boldsymbol F_{J+1}^{+}(s;c)
 +\boldsymbol F_{J-1}^{+}(s;c)
 &=M_+\boldsymbol F_J^{+}(s;c),
 \label{eq:fixed-plus}\\
 \boldsymbol F_{J+1}^{-}(s;c)
 +\boldsymbol F_{J-1}^{-}(s;c)
 &=M_-\boldsymbol F_J^{-}(s;c).
 \label{eq:fixed-minus}
\end{align}
\end{theorem}

\begin{convention}[Formal rank-zero seed]
The $J=1$ instance uses
$\boldsymbol b_0^+=\boldsymbol b_0^-=\boldsymbol m_{n,n}$ as an algebraic
seed.  A relation among three geometric tensor-curl ranks begins at
$J\geq2$.
\end{convention}

\begin{proof}
For $\Re s>3$, the series converge absolutely and
\eqref{eq:fixed-plus}--\eqref{eq:fixed-minus} follow term by term from
\eqref{eq:plus-rank}--\eqref{eq:minus-rank}.  Each side has the
meromorphic continuation of Proposition~\ref{prop:finite-hurwitz}.
Uniqueness of meromorphic continuation proves the identities on
$\C$.
\end{proof}

\subsection{The moving geometric origin}

Let $S$ and $L$ denote the unilateral right and left sequence shifts:
\begin{equation}
 (Sa)(0)=0,\qquad (Sa)(n)=a(n-1)\ (n\geq1),
 \qquad
 (La)(n)=a(n+1).
 \label{eq:unilateral-shifts}
\end{equation}
The symbol $L$ is notation for the left shift; it is not a two-sided
inverse to $S$ because the sequences begin at $n=0$.

\begin{lemma}[Offset reindexing]
\label{lem:offset}
For every quasipolynomial sequence $a$ and $c>1$,
\begin{align}
 F_a(s;c+1)&=F_{Sa}(s;c),
 \label{eq:shift-up}\\
 F_a(s;c-1)&=F_{La}(s;c)+a(0)(c-1)^{-s},
 \label{eq:shift-down}
\end{align}
first in the half-plane of absolute convergence and hence
meromorphically on $\C$.
\end{lemma}

\begin{proof}
Equation~\eqref{eq:shift-up} follows by replacing $n$ with $n+1$.
For \eqref{eq:shift-down}, separate the $n=0$ term and replace
$n\geq1$ by $n-1$.  Proposition~\ref{prop:finite-hurwitz} completes
the continuation.
\end{proof}

\begin{theorem}[Geometric-offset boundary correction]
\label{thm:physical-offset}
Let $J\geq1$, let $c_J=J+1$, and let $M_\pm$ be as in
\eqref{eq:rank-matrices}.  For each branch, define
\[
 \boldsymbol{\mathcal B}_J^\pm(s)
 =
 \boldsymbol F_{(S-I)\boldsymbol b_{J+1}^\pm}(s;J+1)
 +\boldsymbol F_{(L-I)\boldsymbol b_{J-1}^\pm}(s;J+1)
 +\boldsymbol b_{J-1}^\pm(0)J^{-s}.
\]
Then, as vector-valued meromorphic functions on $\C$,
\begin{equation}
 \boxed{\;
 \boldsymbol F_{J+1}^\pm(s;c_{J+1})
 +\boldsymbol F_{J-1}^\pm(s;c_{J-1})
 -M_\pm\boldsymbol F_J^\pm(s;c_J)
 =
 \boldsymbol{\mathcal B}_J^\pm(s).
 \;}
\label{eq:physical-correction}
\end{equation}
At $J=1$, the $J-1$ term is the formal algebraic seed; all three
terms correspond to geometric tensor-curl ranks when $J\geq2$.  In this
boundary case the lowest-mode factor is $1^{-s}=1$, as intended.
\end{theorem}

\begin{proof}
Apply Theorem~\ref{thm:all-s} at the common offset $c=J+1$.
Lemma~\ref{lem:offset} turns the $J+1$ geometric offset into
$S\boldsymbol b_{J+1}^\pm$ and the $J-1$ geometric offset into
$L\boldsymbol b_{J-1}^\pm$ plus the displayed lowest-mode term.
Subtracting the common-offset identity gives
\eqref{eq:physical-correction}.
\end{proof}

\begin{remark}[What the correction means]
\label{rem:cocycle-language}
The right-hand side of \eqref{eq:physical-correction} is an exact
inhomogeneous telescoping correction, or geometric-offset boundary
cocycle.  It measures the cost of changing rank and spectral origin
simultaneously.  The term
$\boldsymbol b_{J-1}^\pm(0)J^{-s}$ is the mode encountered at the
unilateral boundary.  The word ``cocycle'' is descriptive here;
\eqref{eq:physical-correction} does not assert a group-cohomology class.
\end{remark}

\begin{corollary}[Signed and unsigned assembly]
\label{cor:signed-unsigned}
The full eta and absolute-zeta vectors are obtained from
\eqref{eq:physical-correction} by taking the branch difference and
sum, respectively.  The displayed result is the coupled two-matrix
identity; a reduction to one matrix would require additional
invariant-subspace structure not asserted here.
\end{corollary}

\section{Relative zeta functions and log-determinants}

Let $\boldsymbol d=(d_0,\ldots,d_r)^T$ be the dimension vector and let
$\boldsymbol e_0^T\boldsymbol z=z_0$.  Define
\begin{equation}
 R=I-\boldsymbol d\,\boldsymbol e_0^T,
 \qquad
 (R\boldsymbol z)_i=z_i-d_i z_0.
 \label{eq:relative-projection}
\end{equation}
Thus $R^2=R$ and $R\boldsymbol d=0$.  For each column $j$ of either
McKay matrix, taking dimensions in
$W_{L/R}\otimes\rho_j$ gives
$\sum_i d_i(A_{L/R})_{ij}=2d_j$.  The matrices are symmetric, hence
$A_L\boldsymbol d=A_R\boldsymbol d=2\boldsymbol d$, and both rank
matrices satisfy $M_\pm\boldsymbol d=2\boldsymbol d$.

\begin{proposition}[Induced relative recurrence]
\label{prop:relative}
On $\operatorname{im}R$, define
\[
 \widehat M_\pm
 =RM_\pm R|_{\operatorname{im}R}.
\]
For fixed $c>0$, the relative branch vector
\[
 \widehat{\boldsymbol F}_J^\pm(s;c)
 =R\boldsymbol F_J^\pm(s;c)
\]
satisfies
\begin{equation}
 \widehat{\boldsymbol F}_{J+1}^\pm(s;c)
 +\widehat{\boldsymbol F}_{J-1}^\pm(s;c)
 =
 \widehat M_\pm
 \widehat{\boldsymbol F}_J^\pm(s;c)
 \label{eq:relative-recurrence}
\end{equation}
meromorphically on $\C$.
\end{proposition}

\begin{proof}
Apply $R$ to Theorem~\ref{thm:all-s}.  Every vector decomposes as
$\boldsymbol z=R\boldsymbol z+\boldsymbol d z_0$.  Since
$RM_\pm\boldsymbol d=0$,
\[
 RM_\pm\boldsymbol z
 =RM_\pm R\boldsymbol z,
\]
which proves \eqref{eq:relative-recurrence}.
\end{proof}

\begin{proposition}[Pole-free relative rows]
\label{prop:relative-pole-free}
Suppose $\Gamma$ acts freely and effectively on $S^3$.  Then every
component of $R\boldsymbol F_J^\pm(s;J+1)$ is entire.  Consequently
the relative branch log-determinant requires no pole subtraction at any
stage of its continuation.
\end{proposition}

\begin{proof}
By Theorem~\ref{thm:eta-entire}, the only pole-producing class of an
effective free action is the identity.  Its character-projection
weight in coordinate $i$ is $d_i$ times its weight in the trivial
coordinate.  The projection $R$ subtracts exactly that contribution.
All residues therefore vanish, and Proposition~\ref{prop:finite-hurwitz}
leaves no other singularities.
\end{proof}

\begin{corollary}[Relative branch log-determinant]
\label{cor:relative-det}
The functions in \eqref{eq:relative-recurrence} are regular at
$s=0$.  Define the componentwise relative branch log-determinant by
\begin{equation}
 \boldsymbol{\mathcal L}_J^\pm(c)
 =-\partial_s
 \widehat{\boldsymbol F}_J^\pm(s;c)\big|_{s=0}.
 \label{eq:relative-det}
\end{equation}
Then
\begin{equation}
 \boldsymbol{\mathcal L}_{J+1}^\pm(c)
 +\boldsymbol{\mathcal L}_{J-1}^\pm(c)
 =
 \widehat M_\pm\boldsymbol{\mathcal L}_J^\pm(c).
 \label{eq:relative-det-rec}
\end{equation}
At the geometric offsets, differentiating
$R\eqref{eq:physical-correction}$ at $s=0$ gives the corresponding
inhomogeneous relative log-determinant identity.
\end{corollary}

\begin{proof}
The possible poles in Proposition~\ref{prop:finite-hurwitz} are
$1,2,3$, so differentiation at zero is legitimate.  Differentiate
\eqref{eq:relative-recurrence}; the geometric-offset statement follows
in the same way from Theorem~\ref{thm:physical-offset}.
\end{proof}

\begin{proposition}[Exact derivative form]
\label{prop:closed-logdet}
Let a component of a relative branch functional be written in the
finite-Hurwitz form
\begin{equation}
 \widehat F(s)=
 \sum_{r=0}^{P-1}\sum_{h=0}^{2}
 \widehat q_{r,h}P^{h-s}
 \zH\!\left(s-h,x_r\right)
 +\sum_w c_w w^{-s},
 \qquad x_r=\frac{r+c}{P},
 \label{eq:relative-atom-form}
\end{equation}
where the finite point terms record any reindexing used to place the
Hurwitz arguments in a chosen fundamental interval.  Then
\begin{equation}
 -\widehat F'(0)
 =\widehat F_z(0)\log P
  -\sum_{r,h}\widehat q_{r,h}P^h
       \zH'(-h,x_r)
  +\sum_w c_w\log w,
 \label{eq:closed-logdet}
\end{equation}
where $\widehat F_z(0)$ denotes the Hurwitz-zeta part of
$\widehat F(0)$.  Thus every relative determinant row has a finite
exact symbolic decomposition into rational coefficients,
Hurwitz-zeta derivatives at rational arguments, and logarithms.
\end{proposition}

\begin{proof}
Differentiate each finite atom of
\eqref{eq:relative-atom-form} at zero.  The derivative of
$P^{h-s}$ supplies the $\log P$ term, while
$-\partial_s(c_ww^{-s})|_{s=0}=c_w\log w$.
\end{proof}

For a concrete example, take $\Gamma=Q_8\times\Z_3$, with $Q_8$ in the
left slot and $\Z_3$ in the right slot.  Let $E$ denote the
two-dimensional irreducible representation of $Q_8$, and let
$c_0,c_1,c_2$ denote the characters of $\Z_3$.  On the plus branch at
$J=1$, the quasipolynomial period is $P=12$.  Two representative relative
rows are
\begin{center}
\begin{tabular}{lccc}
\toprule
twist & $[\log P](-\widehat F_i'(0))$ & $\eta_{1,i}(0)$ & $-\widehat F_i'(0)$ \\
\midrule
$E\boxtimes c_0$ & $-3/4$ & $-4/9$ & $-0.89893954730931921\ldots$ \\
$E\boxtimes c_1$ & $-5/12$ & $2/9$ & $-0.96202650840501164\ldots$ \\
\bottomrule
\end{tabular}
\end{center}
The eta column is the ordinary signed eta invariant, not the relative
invariant.  The determinant columns are relative.  The decimal column is a display of the exact decomposition
\eqref{eq:closed-logdet}, not its definition.

\begin{remark}[Log-determinant scope]
For each component, the exponential of
\eqref{eq:relative-det} is the zeta-regularized determinant ratio of
that branch to $d_i$ copies of the trivial branch.  The additive
recurrence \eqref{eq:relative-det-rec} is a log-determinant identity.
For arbitrary fixed $c$, it describes auxiliary shifted spectral
products; it is the geometric branch product at rank $J$ only when
$c=c_J=J+1$.  The moving geometric family is governed instead by the
inhomogeneous identity above.  Negative-branch products use absolute
eigenvalue magnitudes and contain no spectral-cut phase.  These formulas
concern the round branch operators; no nonround operator determinant is
asserted.
\end{remark}

\section{Affine-ADE reductions and two-sided examples}

\subsection{One-sided reduction}

Suppose $\Gamma=G\times\{1\}$ with $G< SU(2)$.  If $N$ is the
ordinary McKay adjacency of $G$, then
\[
 A_L=N,\qquad A_R=2I,\qquad
 \boldsymbol m_{p,q}=(q+1)\boldsymbol m_{p,0}.
\]
Consequently,
\[
 M_+=N^2-2I,\qquad M_-=2I.
\]
Thus one TT branch obeys the familiar affine-ADE rank recurrence
while the other has vanishing second difference.  Both are shadows
of the same two-sided theorem.  This also explains why the
one-sided recurrence technology of
\cite{Kostant1985,Steinberg1985,Springer1987,Stekolshchik2008,
Suter2007,Benkart2016,JingWangZhang2020,JingWangZhang2021}
is a direct antecedent but not the full bivariate statement.

\subsection{Product reduction}

For $\Gamma=G_L\times G_R$,
\[
 A_L=N_L\otimes I,
 \qquad
 A_R=I\otimes N_R,
\]
and, after arranging the multiplicities as a matrix,
\begin{align*}
 \mathsf M_{p+1,q}+\mathsf M_{p-1,q}
 &=N_L\mathsf M_{p,q},\\
 \mathsf M_{p,q+1}+\mathsf M_{p,q-1}
 &=\mathsf M_{p,q}N_R^T.
\end{align*}
If $G_L$ and $G_R$ are binary polyhedral, the factors are affine-ADE
adjacencies.  A general nonproduct subgroup has no such outer-product
matrix arrangement; the vector laws of
Theorem~\ref{thm:bivariate} remain the correct formulation.

\subsection{Computational verification}

We checked the formulas independently in exact character and
Hurwitz-atom arithmetic for the mixed products
$Q_8\times\Z_3$, $2T\times\Z_5$, and $Q_8\times\Z_5$, together with
nonproduct cyclic lifts.  The computations recovered the two branch
recurrences, Chebyshev reconstruction, dimension identity
\[
 \sum_i d_i m_{p,q}^i=(p+1)(q+1),
\]
orientation exchange, conjugate-twist symmetry, and the finite-Hurwitz
decompositions.  A deliberately nonfree diagonal $\Z_3$ action exhibited
the resonant residue \eqref{eq:resonant-residue} while its signed branch
difference remained entire.  Direct traces agreed with the heat expansion
inside its convergence disk, and the relative determinant examples showed
exact pole cancellation and the differentiated geometric-offset identity.
These computations test conventions and exceptional cases; the proofs do
not depend on their numerical range.

\section{Relation to prior work}

The representation-theoretic ingredients have several close
antecedents.  McKay's observation and the work of Kostant organize the
one-sided symmetric-power restrictions by affine-ADE data
\cite{McKay1980,Kostant1985,Steinberg1985,Springer1987,
Stekolshchik2008}.  Poincar\'e-series, walk, Cartan, and
Chebyshev formulations were developed further in
\cite{Rossmann2003,Suter2007,Benkart2016,JingWangZhang2020,
JingWangZhang2021}.  Xiong's edge-coloured construction is the nearest
antecedent to the two-operator pair \cite{Xiong2020}; it does not state
the bivariate copy-multiplicity recurrence, Chebyshev grid, TT-line
system, or a spectral application.

For spherical space forms, Ikeda and successors developed
character-theoretic spectral multiplicities and generating functions
\cite{Ikeda1980I,Ikeda1980II,Ikeda1980Lens,
IkedaYamamoto1979,Bar1996,LauretMiatelloRossetti2021}.  Twisted eta
formulas, complete finite-$SU(2)$ Dirac tables, and Molien encodings are
classical \cite{APS1975,Donnelly1978,Gilkey1984,
CisnerosMolina2001,Degeratu2001,Degeratu2009}.  Determinants of
spherical space forms from Dirichlet-series data are also established
\cite{BauerFurutani2008}.  To the best of our knowledge, these sources do
not state the bivariate copy-multiplicity recurrence with its tensor-curl
line system, or the branchwise all-$s$ identity with a moving geometric
offset.  The monographs of du Val and Conway--Smith
\cite{duVal1964,ConwaySmith2003} concern the classification of finite
rotation groups rather than this multiplicity machinery.

\section{Discussion}

The results separate three layers that are easily conflated.

\begin{enumerate}
\item The representation ring determines how sphere harmonics decompose
across topology and flat twist.  The finite matrices $A_L,A_R$ are a
finite exact encoding of this multiplicity layer.
\item The tensor-curl dispersion relation selects the two lines
$q-p=\pm2J$ and assigns the rank-dependent eigenvalue $n+J+1$.
\item Hurwitz continuation converts the quasipolynomial line data into
meromorphic spectral functions and relative branch log-determinants.
\end{enumerate}

The pole theorem adds a fourth structural distinction.  The unsigned
branch sum retains the central polynomial density, hence the local
heat singularities, whereas the signed branch difference discards the
entire pole part and retains only regular spectral asymmetry.  This is
not just regularity at $s=0$: eta-entirety holds throughout the complex
plane, including in nonfree representation-ring examples where
additional resonant poles appear in each branch separately.  For free
effective quotients, the relative projection removes the remaining
identity density and turns every branch row into an entire function.
Thus local heat data, global spectral asymmetry, and relative
determinants are three different finite images of the same bivariate
grid.

At a common offset, the second and third layers preserve the
representation-ring recurrence exactly.  At the geometric offset, the
lowest-mode boundary records the fact that the spectral coordinate
changes with rank.  This is why the same finite algebra can propagate
full spectral functions without erasing their geometric rank dependence.

Several extensions are natural.  The bivariate generating function can
be used to compare inverse-spectral information carried by the
left/right operator pair.  In particular, one may ask whether a finite
rectangle of $\boldsymbol m_{p,q}$, or finitely many TT diagonals,
determines the embedded representation ring up to orientation exchange;
the present reconstruction reduces that question to finite matrix data
but does not solve the resulting equivalence problem.  The relative
log-determinant recurrence gives an exact route to further
round-space-form tables and constrains coincidences across rank: any
putative equality must propagate through both branch matrices and the
explicit boundary term.  A separate
nonround problem is to replace the round representation decomposition
by a mixed-order projected pseudodifferential analysis on Berger
backgrounds.  None of those extensions is needed for the present
theorems, and the last is not implied by them.

\section{Conclusion}

The restrictions of $V_p\boxtimes V_q$ to a finite subgroup of
$\operatorname{Spin}(4)$ form a two-dimensional recurrence governed by
two commuting representation-ring operators.  Its restriction to the
round tensor-curl helicity lines yields two branchwise rank laws and a
diagonal rank--level identity.  Finite-group character theory makes
each line quasipolynomial, so the complete branch spectral function is
a finite Hurwitz-zeta object.  The rank recurrence then holds
meromorphically for all $s$ at fixed offset, while the geometric
rank-dependent origin contributes an explicit unilateral boundary
correction.  Relative zeta functions and determinant ratios inherit the
same structure at the log-determinant level.

In addition, the signed eta function is entire for every finite
two-sided subgroup, while freeness localizes all unsigned poles to the
central polynomial density.  At the geometric offset this density has
the universal form $(n+J+1)^2-J^2$, forcing
$\Res_2=0$ and $\Res_1=-J^2\Res_3$.  These residues generate a
convergent heat expansion, and their cancellation under the relative
projection gives finite exact formulas for pole-free determinant
ratios.  The recurrence therefore propagates not only multiplicities
but also pole data, heat coefficients, eta functions, and
log-determinants, with the same lowest-mode correction whenever the
geometric origin moves.

The affine-ADE theory is therefore not replaced but embedded: it is the
one-sided or product-factor face of a broader two-sided
$\operatorname{Spin}(4)$ representation-ring system.  This viewpoint
turns inverse questions about round tensor-curl spectra into finite problems
for a commuting operator pair and its distinguished diagonals.  Whether a
finite portion of that data determines the embedded representation ring,
and how the recurrence deforms away from the round metric, remain open.

\enlargethispage{\baselineskip}
\section*{Acknowledgment of assistance}
AI-assisted research, computation, checking, and editorial support used
OpenAI Codex and Anthropic Claude. Luqman Mahmood directed the work
and is responsible for the claims and final text. These internal
checks are not external peer review.

\bibliographystyle{alpha}
\bibliography{references}

\end{document}
